\section{Experimental Evaluation}
\label{Sec:experiments}
Our empirical study is designed to test the mechanism suggested by the
theoretical analysis.  Splitting a function-node does not change the objective
function, but it changes the message-passing dynamics by modifying the local
active-minimizer geometry and by introducing short feedback loops between the
copies of each split function-node.  These feedback loops can separate
candidate assignments quickly, but, without damping, they may also generate
oscillations.  Damping is therefore expected to stabilize the oscillatory
component created by splitting.
The experiments are divided into four parts.  First, we compare standard
Min-sum (MS), damped Min-sum (DMS), Min-sum on split constraint factor
graphs (MS-SCFG), and DMS on split constraint factor graphs (DMS-SCFG),
in final cost and in how quickly the selected assignment settles.
Second, we vary the split itself: a random split ratio, a fixed asymmetric
ratio, a temporary asymmetric ratio, and delayed splitting, in which DMS
is performed on the original factor graph for a fixed number of iterations
before the function-nodes are split.  Third, we use the two assignments
produced by the split-only oscillation as a diagnostic of the two-solution
structure of Section~\ref{sec:tworoutes}.  Fourth, we compare against Deep
Attentive Belief Propagation (DABP) and two ablations: DABP without
splitting and DABP with a symmetric split.
\subsection{Experimental Setup}
We evaluated the algorithms on five binary benchmark families, generating 50
independent instances per family (seeds $0$--$49$).
\emph{Random sparse} and \emph{random dense}: 50 agents, one variable each
with $|D_i| = 10$, Erd\H{o}s--R\'enyi constraint graphs with density $p_1 = 0.1$
(sparse) and $p_1 = 0.6$ (dense), and integer constraint costs drawn uniformly
from $[100, 200)$; disconnected components are force-connected.
\emph{Graph coloring}: 50 agents, 3 colors, $p_1 = 0.1$, with not-equal
constraints costing 10 when neighboring agents select the same color and 0
otherwise.
\emph{Scale-free}: Barab\'asi--Albert graphs following \citet{CohenGZ20}: 7
initial randomly connected agents, each new agent attaching preferentially to 3
existing agents, $n = 50$, $|D_i| = 10$, integer costs uniform in $[100, 200)$.
\emph{Meeting scheduling}: following \citet{CohenGZ20}, 90 agents schedule 20
meetings into 20 time slots; each agent participates in two random meetings;
each constrained meeting pair is assigned a travel time drawn uniformly from
$\{6, \ldots, 10\}$, and the cost counts the overbooked (shared) agents when
the slot difference is below the travel time.
All instances add tiny random unary value preferences (uniform in
$[0, 10^{-2})$) for tie breaking, as in prior Max-sum studies
\cite{FarinelliRPJ08}; their total mass stays below the smallest structural
cost gap, so they cannot change which assignment is optimal.
All damped variants use $\lambda = 0.9$.

All algorithms were executed for 2000 synchronous iterations.  At each
iteration, the assignment selected by the current beliefs was evaluated using
the original cost tables.  Thus, even when the algorithm runs on a split
constraint factor graph, all costs are measured with respect to the original
problem.  We report means over the 50 instances per family; standard
deviations accompany all summary statistics, and all pairwise comparisons are
evaluated with paired t-tests and Wilcoxon signed-rank tests over the common
instances.  Full per-benchmark summary and significance tables are included
with the code.
We focus on final solution cost and on the cost trajectory over time.  This is
different from the anytime emphasis of previous work.  The goal here is not
only to determine whether an algorithm visits a good assignment, but whether
the combination of splitting and damping stabilizes to a high-quality assignment
by the end of the run.
The DABP curves are plotted on the same iteration axis after a wall-clock
normalization.  In the plot legend, a label of the form \texttt{x6.3} means that
one iteration of that DABP variant took 6.3 times as long as one DMS iteration
on the same machine.  All such timing measurements were taken on the same
GPU workstation with an NVIDIA RTX 4090 24GB GPU.  Therefore, a DABP curve
labelled \texttt{x6.3} should be read as a wall-clock-adjusted curve, not as a
plain per-iteration curve.
\subsection{Algorithms}
We compare the following algorithms; the labels in parentheses are the ones
used in the plot legends.
\textbf{MS.}  Standard synchronous Min-sum on the original factor graph.
\textbf{DMS.}  Damped Min-sum on the original factor graph.
\textbf{MS-SCFG} (\texttt{MS s=.5})\textbf{.}  Standard Min-sum on a symmetric
split constraint factor graph.  Each original function-node is replaced by two
function-nodes whose cost tables each contain half of the original cost table.
\textbf{DMS-SCFG} (\texttt{DMS s=.5})\textbf{.}  Damped Min-sum on the
symmetric split constraint factor graph.
\textbf{DMS-SCFG-rand} (\texttt{DMS s=.4-.6})\textbf{.}  Damped Min-sum on a
\emph{random} SCFG: each cost-table entry $c$ is split into $u \cdot c$ and
$(1-u) \cdot c$ with $u \sim U[0.4, 0.6)$, the best-performing split reported
by \citet{CohenGZ20}.
\textbf{DMS-SCFG-0.95} (\texttt{DMS s=.95})\textbf{.}  Damped Min-sum on a
fixed asymmetric SCFG: the first copy of every function-node holds $0.95$
of each entry and the second $0.05$ --- the ratio used inside DABP, without
the network.
\textbf{DMS-SCFG-pulse} (\texttt{DMS pulse})\textbf{.}  Damped Min-sum on the
symmetric SCFG with a temporary asymmetric split: before iteration $64$ the
two copies are re-weighted to $0.95/0.05$, and before iteration $256$ back
to $0.5/0.5$; all messages are kept at both changes.  The schedule was fixed
on separate instances (seeds $6000$--$6031$ of the random sparse family) before
these runs.
\textbf{DMS-SCFG-at-$K$} (\texttt{DMS @$K$})\textbf{.}  Delayed splitting.  The
algorithm performs DMS on the original factor graph for $K$ iterations and then
converts the graph to an SCFG, redistributing the previously sent messages
across the two copies.  We evaluate $K\in\{50,100,300,500,1000\}$, and for
random dense graphs also $K=1500$.  Table~\ref{tab:final_costs} reports the
fixed choice $K=1000$ on every benchmark; the per-benchmark choice of $K$ is
evaluated with held-out instances in Section~\ref{sec:heldout}.
\textbf{MS-SCFG-MGM-200 and MS-SCFG-opt-200} (\texttt{MGM@200},
\texttt{Opt merge@200})\textbf{.}  The diagnostic of
Section~\ref{Sec:MS-split}: the two assignments selected at iterations 198
and 199 of the MS-SCFG run --- the two branches of its period-2 oscillation
--- define a two-value menu per variable.  The MGM version applies an
MGM-style local improvement restricted to the menus, started from both
branches, keeping the better result.  The optimal version solves the
restricted binary-menu problem with branch and bound (300-second cap); the
search is exact on every benchmark except random dense, where the reported
value is the best assignment found within the cap, warm-started from the
MGM result.  The inverted variant (\texttt{MGM inv@200}) flips every
disagreement variable of the MGM result to the other value of its menu.
\textbf{Optimal} (\texttt{Opt n=$k$})\textbf{.}  Branch and bound on the
original problem with a 60-second limit per instance, reported only for the $k$
instances on which the search completed (21 of 50 on graph coloring; the search
does not finish on the domain-10 benchmarks).
\textbf{DABP.}  Deep Attentive Belief Propagation~\cite{DengKL022}.  DABP
learns dynamic damping factors and neighbor weights for composing BP messages.
The original DABP implementation also uses an SCFG representation with a
$0.95/0.05$ split ratio.
\textbf{DABP-NoSplit.}  The same attentive architecture, but run on the original
factor graph without splitting.  This ablation tests whether learned damping
and learned neighbor weights are sufficient without the split representation.
\textbf{DABP-SymSplit.}  The same attentive architecture, but using a symmetric
$0.5/0.5$ split instead of the $0.95/0.05$ split.  This ablation tests whether
the advantage of DABP depends on the highly asymmetric split or only on the
presence of a split representation.
\subsection{Final Solution Cost and Settling}
Table~\ref{tab:final_costs} reports the mean final solution cost after 2000
iterations.  Lower values indicate better solutions.
\begin{table*}[t]
\centering
\caption{Mean final solution cost after 2000 iterations (50 instances per
benchmark; lower is better; best in bold).  \texttt{s=.95}: fixed
$0.95/0.05$ split; \texttt{pulse}: temporary $0.95/0.05$ split over
iterations $64$--$255$; \texttt{@1000}: split at iteration $1000$.
Standard deviations and full pairwise significance tables accompany the
code.}
\label{tab:final_costs}
\scriptsize
\setlength{\tabcolsep}{3.5pt}
\begin{tabular}{lrrrrrrrrrr}
\toprule
\textbf{Problem}
& \textbf{MS}
& \textbf{DMS}
& \textbf{MS-SCFG}
& \textbf{DMS-SCFG}
& \textbf{DMS-SCFG s=.95}
& \textbf{DMS-SCFG pulse}
& \textbf{DMS-SCFG @1000}
& \textbf{DABP}
& \textbf{DABP-NoSplit}
& \textbf{DABP-SymSplit} \\
\midrule
Graph coloring
& $666.19$
& $138.81$
& $668.61$
& $34.42$
& \todofill{GC s=.95}
& \todofill{GC pulse}
& $\mathbf{31.63}$
& $39.24$
& $136.41$
& $35.02$ \\
Meeting scheduling
& $87.77$
& $20.70$
& $81.80$
& $8.10$
& \todofill{MS s=.95}
& \todofill{MS pulse}
& $\mathbf{8.04}$
& $27.22$
& $24.00$
& $30.15$ \\
Random dense
& $107935.73$
& $99557.33$
& $107196.39$
& $99758.48$
& \todofill{RD s=.95}
& \todofill{RD pulse}
& $99445.15$
& $\mathbf{99223.68}$
& $100309.79$
& $99665.79$ \\
Random sparse
& $17204.93$
& $15100.29$
& $17073.27$
& $14633.43$
& \todofill{RS s=.95}
& \todofill{RS pulse}
& $14554.99$
& $\mathbf{14519.01}$
& $15108.47$
& $14631.11$ \\
Scale-free
& $19630.55$
& $17079.39$
& $19291.73$
& $16761.71$
& \todofill{SF s=.95}
& \todofill{SF pulse}
& $16636.74$
& $\mathbf{16628.46}$
& $17088.25$
& $16771.71$ \\
\bottomrule
\end{tabular}
\end{table*}
Splitting alone is not sufficient.  MS-SCFG remains close to MS on every
benchmark, and its selected assignment rarely settles: the cost is
constant over the final $100$ iterations in at most $6$ of $50$ runs on
any benchmark, and in none on the random families.  This
is the oscillation analyzed in Sections~\ref{Sec:MS-split}
and~\ref{sec:tworoutes}.
Damping alone improves MS substantially, but on the structured benchmarks
it does not approach the split-damped variants.  In graph coloring, DMS
reaches a mean final cost of $138.81$, while DMS-SCFG reaches $34.42$ and
the delayed split $31.63$.  In meeting scheduling, DMS reaches $20.70$,
while DMS-SCFG reaches $8.10$ and the delayed split $8.04$.  These gaps are
consistent with the claim that splitting changes the local message
geometry in a way that damping alone does not.
On the three random families the picture is different, and the difference
is instructive.  DMS reaches final costs within $3\%$ of the split
variants, and on random dense graphs its final cost is even lower than
that of DMS-SCFG ($99557$ versus $99758$).  What splitting changes there is
not the quality of the assignment but \emph{when} it is reached.
Table~\ref{tab:settling} reports, per algorithm, how many of the $50$ runs
end with a constant cost over the final $100$ iterations and the median
iteration of the last cost change among those runs.  DMS-SCFG settles in
$49$--$50$ of $50$ runs on the random families, at a median iteration of
$43$--$70$; DMS settles in $20$--$40$ runs, at a median iteration of
$359$--$720$, and in $2$ and $8$ runs on meeting scheduling and graph
coloring.  Delayed splitting settles within about $100$ iterations of the split for
every $K$ (median last change at iteration $124$--$174$ for $K=100$ and
$335$--$391$ for $K=300$), or earlier when DMS had already settled.  DABP settles in most runs but late on the dense
and scale-free graphs (median above $1000$), and early on random sparse
graphs.
\begin{table}[t]
\centering
\caption{Settling: runs (of $50$) whose cost is constant over the final
$100$ iterations, and the median iteration of the last cost change among
them.  MS settles in at most $2$ and MS-SCFG in at most $6$ runs per benchmark.}
\label{tab:settling}
\scriptsize
\setlength{\tabcolsep}{3pt}
\begin{tabular}{lrrrrr}
\toprule
& \textbf{GC} & \textbf{MS} & \textbf{Dense} & \textbf{Sparse} & \textbf{SF} \\
\midrule
DMS & $8$ / $379$ & $2$ / $1386$ & $40$ / $720$ & $20$ / $606$ & $22$ / $359$ \\
DMS-SCFG & $40$ / $129$ & $44$ / $154$ & $50$ / $61$ & $49$ / $70$ & $49$ / $43$ \\
DMS-SCFG @1000 & $39$ / $1064$ & $45$ / $1082$ & $50$ / $990$ & $50$ / $1023$ & $49$ / $1006$ \\
DABP & $30$ / $110$ & $14$ / $1114$ & $48$ / $1087$ & $46$ / $31$ & $47$ / $1026$ \\
DABP-NoSplit & $6$ / $346$ & $6$ / $447$ & $37$ / $467$ & $21$ / $406$ & $24$ / $199$ \\
\bottomrule
\end{tabular}
\end{table}
The DABP ablation strengthens this conclusion.  On graph coloring,
DABP-NoSplit reaches $136.41$, close to DMS and far from DABP with splitting
($39.24$) and DMS-SCFG ($34.42$).  Thus, learned damping and learned neighbor
weights do not replace splitting on this structured benchmark.  On meeting
scheduling, both DABP and DABP-NoSplit are much worse than DMS-SCFG.  On the
three random families DABP-NoSplit is statistically indistinguishable from
DMS, while DABP with splitting attains the best final costs on random dense
and random sparse graphs and ties the delayed split on scale-free graphs.
The advantage of the split-damped Min-sum variants is therefore concentrated
on the structured benchmarks; on random dense and sparse graphs, DABP's
learned per-edge parameters win by $0.1$--$0.2\%$, at the price of
centralized execution and a substantially higher per-iteration cost.
Overall, the ablation indicates that splitting remains an essential
structural component even when the algorithm learns dynamic damping
factors and neighbor weights.

\paragraph{Statistical significance.}
All comparisons are paired per instance ($n = 50$); we cite Wilcoxon
signed-rank p-values on final costs, and the paired t-tests agree unless
noted.  Every p-value below $0.05$ cited in this section also survives
Holm's correction over the family of comparisons of its benchmark, except
where marked.
Adding splitting to DMS lowers the final cost significantly on graph
coloring ($p < 10^{-9}$), meeting scheduling ($p < 10^{-8}$), and random
sparse graphs ($p < 10^{-3}$); on scale-free graphs the mean is lower but
the difference is not significant ($p = 0.25$), and on random dense graphs
DMS's final cost is significantly \emph{lower} ($p < 10^{-4}$; t-test
$p = 0.24$), even though DMS has not settled in $10$ of the $50$ runs.
The random $0.4$--$0.6$ split improves on the symmetric split on the three
random families ($p \le 0.006$) and is worse on meeting scheduling
($p < 10^{-4}$).
Removing splitting from DABP is significant on graph coloring
($p < 10^{-9}$), random sparse ($p < 10^{-3}$), and scale-free graphs
($p < 0.005$), but not on meeting scheduling ($p = 0.40$); on random dense
the t-test is significant ($p = 0.003$) while the Wilcoxon test is not
($p = 0.80$).  On graph coloring the differences between DABP and the split
variants are not significant ($p = 0.31$ against DMS-SCFG, $p = 0.96$
against the delayed split); on meeting scheduling every split variant is
significantly better than DABP ($p < 10^{-6}$); on random dense and
random sparse graphs DABP is significantly better than the delayed split
($p = 0.0004$ and $p = 0.007$ at $K=1000$; $p = 0.038$ against $K=1500$ on
random dense) \todofill{DABP vs s=.95 and pulse on dense/sparse}; on
scale-free graphs DABP and the delayed split are tied ($p = 0.75$).
\paragraph{Delayed splitting with held-out selection of $K$.}\label{sec:heldout}
With $K=1000$ fixed, delayed splitting improves significantly on
immediate splitting on random dense ($p < 10^{-5}$), random sparse
($p < 10^{-6}$), and scale-free graphs ($p < 10^{-5}$), and ties it on
graph coloring ($p = 0.30$) and meeting scheduling ($p = 0.76$).  Choosing
$K$ per benchmark on seeds $0$--$24$ and evaluating on seeds $25$--$49$
selects $K = 1500$ on random dense (held-out improvement $431$,
$p < 10^{-5}$), $K = 500$ on scale-free ($104$, $p = 0.002$), and $K = 100$
on the other three, where the held-out improvement is not significant
($p \ge 0.25$).  The gain from delaying the split is therefore established
on the two families whose selection was confirmed out of sample, and the
delayed split is the best distributed variant on every family; on random
dense graphs the held-out gap to DABP is $85$ ($p = 0.05$).
\subsection{Cost Trajectories on Structured Benchmarks}
Figure~\ref{Fig:structured_cost_curves} shows the cost trajectories for graph
coloring and meeting scheduling.  These two benchmarks are the most important
for the main paper because they exhibit the clearest separation between the
mechanisms.  In both cases, MS and MS-SCFG remain high or unstable, while
DMS-SCFG and delayed splitting stabilize at much lower costs.  The DABP-NoSplit
curve is also informative: it tests whether neural attention alone can replace
splitting.  The answer is negative on graph coloring, where removing splitting
keeps DABP close to the DMS region rather than the split-based region.
\begin{figure*}[t]
\centering
\includegraphics[width=.49\textwidth]{plots/graph_coloring_cost.pdf}
\includegraphics[width=.49\textwidth]{plots/meeting_scheduling_cost.pdf}
\caption{Mean solution cost per iteration on the structured benchmarks.
Left: graph coloring.  Right: meeting scheduling.  The DABP labels include a
wall-clock factor: for example, \texttt{DABP x6.3} means that this variant is
6.3 times slower per plotted iteration than DMS on the same RTX 4090 24GB
machine.}
\label{Fig:structured_cost_curves}
\end{figure*}
Because several successful algorithms become close after the transient phase,
Figure~\ref{Fig:structured_cost_curves_zoom} shows a zoomed view of the
lower-cost region.  These plots should be read as convergence plots rather than
anytime plots.  They show whether the decoded solution cost remains low near
the end of the run.
\begin{figure*}[t]
\centering
\includegraphics[width=.49\textwidth]{plots/graph_coloring_cost_zoom.pdf}
\includegraphics[width=.49\textwidth]{plots/meeting_scheduling_cost_zoom.pdf}
\caption{Zoomed view of the lower-cost region on graph coloring and meeting
scheduling.  The zoom highlights the separation between DMS-SCFG, delayed
splitting, DABP, DABP-NoSplit, and DABP-SymSplit after the early transient.}
\label{Fig:structured_cost_curves_zoom}
\end{figure*}
\subsection{The Symmetric-Split DABP Ablation}
The DABP-SymSplit line has a specific interpretation.  It keeps the DABP
architecture and online learning procedure fixed, but replaces the $0.95/0.05$
split with a symmetric $0.5/0.5$ split.  Thus it isolates the effect of the split
ratio.
The $0.95/0.05$ split is significantly better than the symmetric split on the
three unstructured families (random sparse, random dense, and scale-free;
Wilcoxon $p < 10^{-8}$ in each case), while on the two structured benchmarks
the two split ratios are statistically indistinguishable (graph coloring
$p = 0.44$, where the symmetric split has the better mean, $35.02$ vs.\
$39.24$; meeting scheduling $p = 0.23$).
The original $0.95/0.05$ split gives one copy almost the entire original cost
table and gives the second copy only a small fraction.  Early in the run, the
dominant copy is therefore likely to select active minimizers similar to those
selected by the original unsplit function-node.  At the same time, the small
copy already creates an additional feedback channel.  Later in the run, this
weak channel can become more important through accumulated messages and
through the learned DABP neighbor weights.  This provides a plausible mechanism
for why the $0.95/0.05$ split yields better final costs on the unstructured
families: it preserves the original factor behavior early, while leaving a
learned secondary route that can be amplified later.  Proposition~\ref{prop:asym}
gives the corresponding formal handle: the split ratio moves the round-trip
constants and thresholds, and hence the absorption condition.
The symmetric split behaves differently.  Both copies immediately operate at
equal scale.  This makes the feedback more balanced from the first iterations
and removes the strong-copy/weak-copy asymmetry.  The symmetric split is
often faster in wall-clock factor on the structured and sparse benchmarks: the
measured per-iteration ratios give graph coloring \texttt{DABP x6.3} versus
\texttt{DABP sym-split x6.1}, meeting scheduling \texttt{x8.2} versus
\texttt{x7.3}, and random sparse \texttt{x6.9} versus \texttt{x4.6}.  The
effect is not universal: on random dense and scale-free graphs the symmetric
split is not faster by this timing factor.  Since the two split DABP variants
have the same number of factor-node copies, this runtime difference should not
be explained by graph size alone.  The likely explanation is that the split
ratio changes the learned message dynamics and the early stopping / convergence
behavior, rather than the raw number of messages.
This ablation suggests a useful future direction.  We should inspect the learned
DABP neighbor weights and $R$-message magnitudes on the two split copies.  If
the interpretation above is correct, then the $0.95/0.05$ split should show a
weak copy that is initially small but later receives non-negligible learned
weight on some edges.  Such an analysis would connect the empirical DABP
advantage directly to the feedback-loop mechanism studied in this paper.
\subsection{Selecting within the Two-Branch Menus}
Section~\ref{Sec:MS-split} read the split-only oscillation as two
solutions visited out of phase, and proposed a diagnostic: restrict every
variable to the two values it alternates between and select within these
menus.  Table~\ref{tab:merge} compares the raw MS-SCFG final cost with the
MGM and exact selections, with the inverted selection, and with DMS-SCFG.
\begin{table*}[t]
\centering
\caption{Selecting within the two-branch menus of the split-only
oscillation (branches at iterations 198 and 199).  Lower is better.  Both
selections improve on raw MS-SCFG on every benchmark (Wilcoxon
$p < 10^{-8}$); the inverted selection is significantly worse than the MGM
selection on every benchmark ($p < 10^{-9}$); and DMS-SCFG is significantly
better than the exact selection on every benchmark ($p < 10^{-4}$).}
\label{tab:merge}
\begin{tabular}{lrrrrr}
\toprule
\textbf{Problem}
& \textbf{MS-SCFG}
& \textbf{MGM selection}
& \textbf{Exact selection}
& \textbf{Inverted selection}
& \textbf{DMS-SCFG} \\
\midrule
Graph coloring
& $668.61$ & $98.19$ & $83.19$ & $235.61$ & $34.42$ \\
Meeting scheduling
& $81.80$ & $20.98$ & $20.14$ & $25.66$ & $8.10$ \\
Random dense
& $107196.39$ & $100419.69$ & $100419.21$ & $103190.69$ & $99758.48$ \\
Random sparse
& $17073.27$ & $14800.23$ & $14761.37$ & $15369.99$ & $14633.43$ \\
Scale-free
& $19291.73$ & $16942.89$ & $16878.09$ & $17449.27$ & $16761.71$ \\
\bottomrule
\end{tabular}
\end{table*}
The selections recover most of the quality that the uncoordinated
oscillation hides: on every benchmark they close between $82\%$ and
$95\%$ of the gap between raw MS-SCFG and DMS-SCFG, although they only
ever see the two branch assignments.  This is what the two-solution
reading predicts, and the inverted selection adds a second prediction: it
is significantly worse than the MGM selection everywhere, yet still far
better than the raw snapshots, because the other choice from the menus is
again a coherent assignment --- on bipartite graphs, the second rephasing
of Corollary~\ref{cor:bipartite}.  The selections are not, however,
competitive solvers.  DMS-SCFG is significantly better than the exact
selection on every benchmark, and plain DMS, which never splits, beats it
on random dense graphs and ties it on meeting scheduling and scale-free
graphs.  The information carried by the two branches
is real but incomplete, and damping, which resolves the alternation inside
the message passing (Section~\ref{sec:whydamping}), uses it better than any
selection made after the fact.
\subsection{Measuring the Mechanism}
\label{sec:mechanism}
The theory of Sections~\ref{sec:setting} and~\ref{sec:tworoutes} predicts
specific, measurable behavior on benchmark instances, which we instrumented
directly (full protocol, tables, and raw results in
Appendix~\ref{app:tworoutes}).  (i)~\emph{Commitment}: the fraction of
messages with a stable active minimizer rises within ten iterations to
$0.73$--$0.97$ on the split random families (versus $0.10$--$0.36$ unsplit
in the tail, family-dependent), while graph coloring commits only in the
weaker selection-content sense ($0.03$ at iteration $9$, $0.36$ in the
tail) --- commitment is real but partial, which is why the reduction of
Section~\ref{sec:tworoutes} is stated conditionally.  (ii)~\emph{Selection
rule}: on an audited subset of $32$ period-2 runs, the two alternating
solutions mapped to one another under the synchronous local selection rule
exactly, and the rule predicted the next assignment correctly on all
sampled transitions.  (iii)~\emph{Attractor structure}: on the same $32$
runs, the limit pair was a layer-wise local minimum of the alternation
cost $\mathrm{cost}_2$ with zero violating single-variable moves, and on
oscillating bipartite instances the rephasings of
Corollary~\ref{cor:bipartite} were coordinatewise local minima of the
original problem in every case.  (iv)~\emph{Damping}: the smallest grid
$\lambda$ from which all larger values stabilized the assignment tail was
at most $0.4$ on every scanned benchmark instance --- far below the
customary $0.9$, and consistent with Theorem~\ref{thm:lambdastar} ---
while Proposition~\ref{prop:kn} exhibits a family where the required
damping exceeds $0.9$ and grows with degree.
\subsection{Ternary Benchmarks}
We also include a ternary benchmark suite in the appendix.  In these benchmarks
the factors are genuine arity-3 constraints rather than collections of binary
constraints.  The ternary suite is important because the cleanest theoretical
results in this paper are for binary domains and binary local minimizer maps.
For domain size 2, preservation of an active minimizer can force a full neutral
step in which the incoming message difference is cancelled.  With three values,
this clean cancellation need not occur for all rows or columns simultaneously.
The ternary experiments therefore test whether the qualitative split-damped
phenomenon persists when the local minimizer structure is richer and only
semi-neutral behavior may remain.
The ternary plots should not be interpreted as a replacement for the binary
theory.  Rather, they are evidence that the mechanism is not limited to a
single pairwise artifact.  If split-damped variants also stabilize on true
arity-3 factors, then the threshold and minimizer-switching explanation is
likely to extend beyond the simplest binary case, although the full theorem
will require additional case distinctions.
\subsection{Discussion}
The empirical results support the mechanism proposed in the theoretical
sections.  Standard Min-sum performs poorly on cyclic benchmark graphs.
Damping alone improves the result but does not fully exploit the structural
change created by splitting.  Splitting alone changes the message dynamics but
often produces oscillatory behavior.  Combining splitting with damping produces
much lower final costs, especially on structured benchmarks such as graph
coloring and meeting scheduling.
The delayed-splitting results show that splitting is not merely a preprocessing
transformation.  The point at which the split is introduced affects the final
result, suggesting that DMS can first form useful beliefs on the original graph
and that splitting then settles them within about a hundred iterations
(Table~\ref{tab:settling}) by changing the local threshold structure of the
function-node updates.  The improvement over immediate splitting is
established out of sample on random dense and scale-free graphs and is a
tie on the other families; selecting $K$ automatically is an open question.
The split ratio matters as well: the random $0.4$--$0.6$ split, the fixed
$0.95/0.05$ split, and the temporary $0.95/0.05$ pulse all move the final
cost on the random families \todofill{one sentence with the s=.95 and pulse
outcome once the 50-instance runs are in}, which is the effect
Proposition~\ref{prop:asym} makes available: the ratio moves the round-trip
constants and thresholds, and with them the absorbing regime.
The DABP comparison further strengthens the main claim.  Even learned dynamic
damping and learned neighbor weights do not remove the need for splitting.  The
no-split attentive ablation fails to reach the split-based performance on the
structured benchmarks and is indistinguishable from DMS on the random
families, and the original DABP design itself uses a split representation.  The symmetric-split DABP ablation shows that the exact split
ratio also matters on the unstructured families, while on the structured
benchmarks the two ratios are statistically indistinguishable and the symmetric
split often runs faster, providing a cleaner mechanistic baseline.
Understanding this tradeoff will require direct analysis of the learned edge
weights and the messages flowing through the two split copies.
\paragraph{Limitations.}
The formal results of Section~\ref{sec:setting} concern a single split binary
function-node with two-value domains; undamped updates, except for the
conditional Proposition~\ref{prop:damped}; and influence from the remainder of
the graph modeled as a bounded unary perturbation.
Proposition~\ref{prop:asym} extends the absorption analysis to asymmetric
splits, but the unconditional convergence theorem
(Theorem~\ref{thm:convergence}) is proved only for the symmetric case.
Whether the bounded-perturbation condition of Theorem~\ref{thm:dyncvg} holds
on benchmark instances is an empirical question, which the instrumentation of
Section~\ref{sec:mechanism} answers: commitment is partial and
family-dependent.  Accordingly, the oscillation analysis of
Sections~\ref{sec:tworoutes}--\ref{sec:whydamping} is conditional: the
reduction to synchronous local selection assumes full commitment; the
two-solution theorem assumes the selection recursion holds along the run
(verified on every audited instance, but not proved to occur); the damping
threshold concerns fixed, fully committed binary patterns; and
Proposition~\ref{prop:kn} shows that no problem-independent damping value
exists, so the reliability of $\lambda = 0.9$ is a property of the
benchmark families.  Finally, on random dense and
random sparse graphs DABP's final costs remain significantly better than
those of the best distributed variant (by $0.1$--$0.2\%$), and on scale-free
graphs the two are tied; our contribution there is the mechanism
explanation, the settling speed, and the distributed-execution
compatibility, not state-of-the-art solution quality.
